% Figure: SIR phase plane in normalised coordinates (s = S/N, i = I/N).
% Trajectories i(s) = i0 + s0 - s + (1/R0) ln(s/s0) for three values of
% the basic reproduction number R0, all starting near (s, i) = (1, 0).
% The vertical dashed lines mark the susceptible thresholds s = 1/R0 at
% which each trajectory peaks. Curves drawn from the closed-form i(s).
\begin{figure}[htbp]
\centering
\begin{tikzpicture}[
  axis/.style={draw=engblack, -{Stealth}, thick},
  curvea/.style={draw=engblue, very thick},
  curveb/.style={draw=engorange, very thick},
  curvec/.style={draw=engred, very thick},
  thresh/.style={draw=enggray, dashed, thin},
  label/.style={font=\small},
  scale=9,
]
% Axes: s (=S/N) from 0 to 1 on x; i (=I/N) from 0 to 0.4 on y (scaled).
% We scale i by 2 vertically so 0.4 maps to 0.8 for a readable aspect.
\draw[axis] (0,0) -- (1.08,0) node[anchor=west, label] {$s = S/N$};
\draw[axis] (0,0) -- (0,0.92) node[anchor=south, label] {$i = I/N$};
\foreach \x in {0.2,0.4,0.6,0.8,1.0} {
  \draw (\x,-0.012) -- (\x,0.012);
  \node[font=\scriptsize, below] at (\x,-0.02) {\x};
}
\foreach \iv/\iy in {0.1/0.2, 0.2/0.4, 0.3/0.6} {
  \draw (-0.012,\iy) -- (0.012,\iy);
  \node[font=\scriptsize, left] at (-0.02,\iy) {\iv};
}
% Trajectory i(s) = i0 + s0 - s + (1/R0) ln(s/s0), s0 = 0.999, i0 = 0.001.
% Vertical axis scaled by 2.0 (so plotted y = 2 * i).
% R0 = 1.5 (blue): peak at s = 1/1.5 = 0.667
\draw[curvea] plot[domain=0.36:0.999, samples=120]
  ({\x}, {2*(0.001 + 0.999 - \x + (1/1.5)*ln(\x/0.999))});
% R0 = 2.5 (orange): peak at s = 0.40
\draw[curveb] plot[domain=0.11:0.999, samples=140]
  ({\x}, {2*(0.001 + 0.999 - \x + (1/2.5)*ln(\x/0.999))});
% R0 = 3.5 (red): peak at s = 1/3.5 = 0.2857
\draw[curvec] plot[domain=0.057:0.999, samples=150]
  ({\x}, {2*(0.001 + 0.999 - \x + (1/3.5)*ln(\x/0.999))});
% Threshold lines at s = 1/R0
\draw[thresh] (0.667,0) -- (0.667,0.18);
\draw[thresh] (0.40,0) -- (0.40,0.55);
\draw[thresh] (0.2857,0) -- (0.2857,0.78);
% Labels
\node[font=\scriptsize, engblue, anchor=west] at (0.70,0.20) {$\mathcal{R}_{0}=1.5$};
\node[font=\scriptsize, engorange, anchor=west] at (0.43,0.57) {$\mathcal{R}_{0}=2.5$};
\node[font=\scriptsize, engred, anchor=west] at (0.31,0.80) {$\mathcal{R}_{0}=3.5$};
% Start point marker
\draw[engblack, fill=engblack] (0.999,0) circle (0.4pt);
\node[font=\scriptsize, anchor=south west] at (0.86,0.005) {start};
\end{tikzpicture}
\caption[SIR phase plane across reproduction numbers]{Phase-plane
trajectories of the SIR model in normalised coordinates, $i = I/N$
against $s = S/N$, for three basic reproduction numbers. Each
trajectory starts at the lower right with almost the whole population
susceptible and a small infectious seed, follows the closed-form curve
$i(s) = i_{0} + s_{0} - s + \mathcal{R}_{0}^{-1}\ln(s/s_{0})$ leftward
as susceptibles are depleted, peaks where $s = 1/\mathcal{R}_{0}$
(dashed lines), and returns to the $s$-axis at a final susceptible
fraction $s_{\infty} > 0$. A larger $\mathcal{R}_{0}$ raises the peak,
moves the threshold to the left, and leaves fewer survivors uninfected.
The flow is one-way: there is no closed orbit, because removed
individuals do not return to the susceptible pool. The vertical axis is
scaled for readability.}
\label{fig:vol02:ch07:sir-phase}
\end{figure}
